\documentclass[11pt,a4paper]{article}

\usepackage[utf8]{inputenc}
\usepackage[T1]{fontenc}
\usepackage{lmodern}
\usepackage[margin=2.4cm]{geometry}
\usepackage{amsmath,amssymb,amsthm,mathtools}
\usepackage{bm}
\usepackage{booktabs,array,float,multirow}
\usepackage{enumitem}
\usepackage{microtype}
\usepackage{syntax}
\usepackage{xcolor}
\usepackage[numbers,sort&compress]{natbib}
\usepackage{listings}
\usepackage[colorlinks=true,linkcolor=blue!45!black,citecolor=blue!45!black,urlcolor=blue!45!black]{hyperref}
\usepackage{cleveref}

\newtheoremstyle{thmstyle}{\topsep}{\topsep}{\itshape}{}{\bfseries}{.}{.5em}{}
\theoremstyle{thmstyle}
\newtheorem{theorem}{Theorem}[section]
\newtheorem{lemma}[theorem]{Lemma}
\newtheorem{proposition}[theorem]{Proposition}
\newtheorem{corollary}[theorem]{Corollary}
\theoremstyle{definition}
\newtheorem{definition}[theorem]{Definition}
\newtheorem{example}[theorem]{Example}
\theoremstyle{remark}
\newtheorem{remark}[theorem]{Remark}

\newcommand{\N}{\mathbb{N}}
\newcommand{\R}{\mathbb{R}}
\newcommand{\flo}{\beta}
\newcommand{\kap}{\kappa}
\newcommand{\sep}{\sigma}
\newcommand{\med}{\mathfrak{M}}
\newcommand{\Env}{\Gamma}
\newcommand{\store}{\Sigma}
\newcommand{\steps}{\longrightarrow}
\newcommand{\abs}[1]{\lvert#1\rvert}
\newcommand{\seek}{\mathbf{seek}}
\newcommand{\excl}{\mathbf{excluding}}
\newcommand{\via}{\mathbf{via}}
\newcommand{\until}{\mathbf{until}}
\newcommand{\Ty}[2]{#1\ \vdash\ #2}

\lstdefinelanguage{mck}{
  morekeywords={seek,excluding,via,until,closure,catalyst,receiver,observable,type,let,report,decline,over,where,as,floor,record},
  sensitive=true,
  morecomment=[l]{\#},
  morestring=[b]",
}
\lstdefinestyle{mckstyle}{
  language=mck,
  basicstyle=\ttfamily\small,
  keywordstyle=\color{blue!70!black}\bfseries,
  commentstyle=\color{green!45!black}\itshape,
  stringstyle=\color{red!55!black},
  frame=single,
  breaklines=true,
  columns=fullflexible,
  showstringspaces=false,
}

\title{\textbf{Mekaneck: A Substrate-Neutral Language\\[2pt]
for Individuation-Structured Inquiry}\\[4pt]
\large Grammar, Type System, and Operational Semantics of \texttt{.mck}}

\author{%
Kundai Farai Sachikonye\\
\small Technical University of Munich\\
\small \texttt{kundai.sachikonye@tum.de}}

\date{\today}

\begin{document}
\maketitle

\begin{abstract}
\noindent
We define \emph{Mekaneck} (source extension \texttt{.mck}), a small language
whose sole computational primitive expresses the individuation of a target
from a substrate by explicit exclusion, and whose termination condition is
the exhaustion of reachable outcomes rather than the crossing of a
confidence threshold. The language has one binding form, one primitive
\texttt{seek} with a \emph{mandatory} exclusion clause, and a fixed set of
substrate declarations that bind an experimental domain to the primitive
through four obligations: what the receivers are, what observable is
recorded, what counts as a discriminating event, and how the irreducible
floor is estimated. We give the complete lexical and context-free grammar,
a type system with five types and two non-standard rules---a
\emph{positivity} rule rejecting programs that assert an attainable zero
residual, and a \emph{coherence} rule rejecting claims grounded on fewer
than three mutually supporting independent sources---and a small-step
operational semantics over a configuration comprising a substrate binding, a
value store, and a monotone committed record. We prove progress and
preservation for the type system, and hence that a well-typed program does
not get stuck; we prove that evaluation is deterministic given the substrate
and the sequence of catalyst outcomes, so that non-determinism enters only
through substrate responses; we prove that the committed record is strictly
monotone, so a program never returns to a previous configuration and its
provenance order is intrinsic; we prove that every well-typed \texttt{seek}
over a finite catalyst registry terminates, and that it terminates in
exactly one of two typed outcomes---a \emph{resolved} cell or an explicit
\emph{decline} carrying the plurality of incompatible cells found. We prove
that the exclusion clause cannot be eliminated: a \texttt{seek} without it
is not merely unidiomatic but untypable, because the target of an
individuation is not determined by a positive description alone. We then
show the language is substrate-neutral by exhibiting three bindings---an
oscillatory-coherence substrate, a discrete-transition substrate, and a
latency substrate---each requiring only the four obligations, and we
identify precisely the condition on a substrate under which its
discriminating events carry no information, so that a binding may be
diagnosed as uninformative rather than the inquiry being reported as
negative.
\end{abstract}

\noindent\textbf{Keywords:} domain-specific language; operational semantics;
type soundness; individuation by exclusion; closure-based termination;
substrate neutrality; scientific workflow languages.

\tableofcontents

% =====================================================================
\section{Introduction}
\label{sec:intro}
% =====================================================================

\subsection{The design question}

Languages for scientific computation generally supply a rich vocabulary of
operations over a domain and leave the structure of inquiry implicit in how
a user composes them. A script computes some quantities, applies a test,
and---by a convention external to the language---stops when a number falls
below a threshold. The structure of the inquiry lives in the analyst's head
and in comments.

Mekaneck inverts this. It supplies almost no domain vocabulary and instead
makes the structure of inquiry the only thing the language expresses. There
is one computational primitive. It states what is being individuated, what
it is being individuated \emph{against}, by what means, and under what
condition the individuation is complete. Everything domain-specific enters
through a substrate binding that satisfies four obligations and is otherwise
opaque to the language.

Two commitments distinguish the design, and both are enforced by the type
system rather than by convention.

\emph{Exclusion is mandatory.} A \texttt{seek} must state what its target is
being told apart from. This is not stylistic. \Cref{thm:exclusion-necessary}
shows that a positive description alone does not determine a target, so a
\texttt{seek} lacking an exclusion clause has no well-defined denotation and
is rejected at compile time.

\emph{Termination is by exhaustion, not by threshold.} A \texttt{seek}
completes when no remaining catalyst can carry it to an outcome outside the
set already reached. \Cref{thm:closure-stronger} shows this condition is not
obtainable by lowering a threshold, and \Cref{thm:dichotomy} shows it
resolves into exactly two typed outcomes, one of which is an explicit
declination.

\subsection{What the language is not}

Mekaneck is not a general-purpose language and has no facilities for
arithmetic beyond what the substrate supplies, no user-defined functions, no
recursion, and no I/O other than through substrate bindings. It is a
coordination language for inquiry, in the sense that a build language is a
coordination language for compilation: the work happens elsewhere, and the
language expresses only its structure. A program that needs to compute a
spectral estimate calls a substrate that computes it; the language sees an
opaque value.

\subsection{Relation to existing work}

The formal apparatus is standard: a context-free grammar in the style of
\citet{aho2006}, a type system presented as inference rules with progress
and preservation in the manner of \citet{wright1994} and
\citet{pierce2002}, and a small-step structural operational semantics
following \citet{plotkin1981,kahn1987}. The novelty is not in the metatheory
but in what the rules enforce.

Scientific workflow languages---Make and its lineage, and the modern
dataflow workflow engines \citep{koster2012,direct2017,amstutz2016}---express
dependency and let the user supply the stopping condition. Statistical
workflow systems embed stopping rules in libraries rather than in the
language. Sequential-analysis stopping rules
\citep{wald1945,lindley1956,chaloner1995} halt when an expected gain or a
likelihood ratio crosses a boundary; \Cref{thm:closure-stronger} distinguishes
our criterion from all such rules. The requirement of three mutually
supporting independent sources (\Cref{sec:coherence}) is related to the
triangulation heuristics discussed in methodological work on convergent
evidence \citep{munafo2018} and to the treatment of conflicting evidence in
belief functions \citep{shafer1976}, but is here a typing rule with a proof
rather than a recommendation. The insistence that the language report a
plurality when evidence conflicts, rather than select a representative,
parallels the ``honest decline'' available in conformal and abstention-based
prediction \citep{chow1970,vovk2005}.

% =====================================================================
\section{The substrate interface}
\label{sec:substrate}
% =====================================================================

Everything domain-specific enters through a substrate. A substrate is
required to discharge exactly four obligations, and the language depends on
nothing else about it.

\begin{definition}[Substrate]
\label{def:substrate}
A \emph{substrate} $\mathcal{D}$ is a four-tuple of obligations:
\begin{enumerate}[label=\textup{(S\arabic*)},leftmargin=3.2em,nosep]
\item \textbf{Receivers.} A finite set $R(\mathcal{D})$ of bounded systems,
each of which is the unit to which a floor attaches.
\item \textbf{Observable.} For each receiver, a finite sequence of
\emph{states}, each carrying a real-valued uncertainty $S$ and a label drawn
from a finite label set.
\item \textbf{Events.} A rule determining which ordered pairs of states
constitute \emph{discriminating events}, and a type for each such event
drawn from a finite set.
\item \textbf{Floor.} An estimator returning $\flo(\mathcal{D},r)>0$ for
each receiver $r$, the irreducible residual below which uncertainty does not
fall.
\end{enumerate}
\end{definition}

\begin{remark}[The floor obligation is the substantive one]
\label{rem:floor-obligation}
(S1)--(S3) are bookkeeping: they say what the units are, what is measured,
and what counts as a step. (S4) is a commitment with empirical content. A
substrate that discharges it by returning the minimum observed $S$ has
discharged it trivially---the value is positive whenever the sample is, and
carries no information. A substrate that discharges it by estimating the
asymptote of separation cost as the sample grows has made a claim capable of
being wrong. The language cannot tell these apart, and we therefore require
a substrate to declare which estimator it uses, so that a reader of a
\texttt{.mck} program can see it. \Cref{sec:diagnostics} gives the
diagnostic that detects the trivial case from data.
\end{remark}

\begin{definition}[Catalyst]
\label{def:catalyst}
A \emph{catalyst} is a named means of advancing an inquiry: a computation
over substrate data, a query to an external source, or an inferential
procedure. Catalysts are opaque to the language, which knows only their
names, their declared independence relation
(\Cref{sec:coherence}), and the outcomes they return.
\end{definition}

\begin{definition}[Cell]
\label{def:cell}
A \emph{cell} is an equivalence class of outcomes under a substrate-supplied
indistinguishability relation. Cells, not points, are the values a
\texttt{seek} returns.
\end{definition}

\begin{remark}[Why the return type is a cell]
\label{rem:cell}
The language cannot return a point-valued answer, because the floor
obligation (S4) asserts a positive irreducible residual and a point-valued
answer would assert a zero residual. This is enforced by the positivity rule
of \Cref{sec:typing} and is the reason the type \texttt{Point} does not
exist in the language.
\end{remark}

% =====================================================================
\section{Lexical structure and grammar}
\label{sec:grammar}
% =====================================================================

\subsection{Token classes}

\begin{definition}[Tokens]
\label{def:tokens}
The lexer recognises: \emph{keywords}
\texttt{substrate}, \texttt{receivers}, \texttt{observable},
\texttt{events}, \texttt{floor}, \texttt{catalyst}, \texttt{independent},
\texttt{let}, \texttt{seek}, \texttt{excluding}, \texttt{via},
\texttt{until}, \texttt{closure}, \texttt{report}, \texttt{resolved},
\texttt{decline}, \texttt{record};
\emph{identifiers} matching \texttt{[A-Za-z\_][A-Za-z0-9\_]*} that are not
keywords; \emph{string literals} delimited by double quotes;
\emph{numeric literals} in decimal notation; \emph{punctuation}
\texttt{\{ \} ( ) , ; : =}; and \emph{comments} from \texttt{\#} to
end-of-line, discarded. Whitespace is insignificant except as a token
separator.
\end{definition}

\subsection{Context-free grammar}

\begin{grammar}
<program> ::= <decl>*

<decl> ::= <substrate-decl> | <catalyst-decl> | <let-decl> | <report-decl>

<substrate-decl> ::= `substrate' <ident> `\{'
  `receivers' `:' <expr> `;'
  `observable' `:' <expr> `;'
  `events' `:' <expr> `;'
  `floor' `:' <expr> `;' `\}'

<catalyst-decl> ::= `catalyst' <ident> `:' <expr>
  [ `independent' <ident-list> ] `;'

<let-decl> ::= `let' <ident> `=' <seek-expr> `;'

<report-decl> ::= `report' <ident> `;'

<seek-expr> ::= `seek' <expr>
  `excluding' <expr>
  [ `via' `(' <ident-list> `)' ]
  `until' <term-cond>

<term-cond> ::= `closure'

<ident-list> ::= <ident> ( `,' <ident> )*

<expr> ::= <ident> | <string> | <number> | <ident> `(' <arg-list>? `)'

<arg-list> ::= <expr> ( `,' <expr> )*
\end{grammar}

\begin{proposition}[The grammar is unambiguous]
\label{prop:unambiguous}
The grammar of \Cref{sec:grammar} is LL(1) and therefore unambiguous.
\end{proposition}

\begin{proof}
Each production for a non-terminal begins with a distinct terminal: the four
\texttt{<decl>} alternatives begin with the distinct keywords
\texttt{substrate}, \texttt{catalyst}, \texttt{let}, \texttt{report}; the
\texttt{<seek-expr>} clauses are introduced by the distinct keywords
\texttt{seek}, \texttt{excluding}, \texttt{via}, \texttt{until}; the
optional clauses \texttt{via} and \texttt{independent} are distinguished
from what follows by their leading keyword; and \texttt{<expr>} alternatives
are distinguished by token class, with the function-call form disambiguated
by a one-token lookahead for \texttt{(}. Hence a single token of lookahead
selects a unique production at every point, so the grammar is LL(1), and
LL(1) grammars are unambiguous \citep{aho2006}.
\end{proof}

\begin{remark}[The exclusion clause is not optional in the grammar]
\label{rem:excl-grammar}
Note that \texttt{via} and \texttt{independent} are bracketed as optional in
the grammar while \texttt{excluding} is not. This is deliberate and is the
syntactic reflection of \Cref{thm:exclusion-necessary}: a \texttt{seek}
without an exclusion clause is not a parse error to be recovered from but a
form the language does not contain.
\end{remark}

\subsection{An example program}

\begin{lstlisting}[style=mckstyle]
# Individuate a coherence regime against the rest of the record.
substrate Osc {
  receivers  : recordings("cohort-A");
  observable : coherence_index();
  events     : label_change();
  floor      : asymptotic_separation();   # not sample_minimum()
}

catalyst spectral  : band_decomposition()  independent surrogate, phase;
catalyst surrogate : phase_randomised()    independent spectral, phase;
catalyst phase     : locking_value()       independent spectral, surrogate;

let regime =
  seek        target_state("high-coherence")
  excluding   all_other_states()
  via         (spectral, surrogate, phase)
  until       closure;

report regime;
\end{lstlisting}

The program declares a substrate discharging the four obligations, three
mutually independent catalysts, and one \texttt{seek} whose target is
individuated against the explicitly named remainder. It terminates by
closure, and \texttt{report} emits either a resolved cell or a declination.

% =====================================================================
\section{The type system}
\label{sec:typing}
% =====================================================================

\subsection{Types}

\begin{definition}[Types]
\label{def:types}
The types are
\[
T \ ::=\ \texttt{Substrate}\ \mid\ \texttt{Catalyst}\ \mid\ \texttt{Cell}\ \mid\ \texttt{Outcome}\ \mid\ \texttt{Floor}.
\]
There is no type of points: the language cannot express a residue-free
answer.
\end{definition}

\begin{definition}[Outcome]
\label{def:outcome}
$\texttt{Outcome}=\texttt{Resolved}(\texttt{Cell})\ \mid\
\texttt{Declined}(\{\texttt{Cell}\})$, a sum whose second summand carries a
set of at least two mutually incompatible cells.
\end{definition}

\subsection{Judgements and standard rules}

A typing environment $\Env$ maps identifiers to types. The judgement
$\Env\vdash e:T$ is derived by the following rules, in which
\textsc{T-Var}, \textsc{T-Sub}, \textsc{T-Cat} are standard.

\begin{align*}
&\textsc{T-Var}\quad
\frac{x:T\in\Env}{\Env\vdash x:T}
\qquad
\textsc{T-Sub}\quad
\frac{\Env\vdash r,o,v,f\ \text{well-formed}}
     {\Env\vdash \texttt{substrate}\,\{r;o;v;f\}:\texttt{Substrate}}\\[6pt]
&\textsc{T-Cat}\quad
\frac{\Env\vdash e:\_\qquad \iota\subseteq\mathrm{dom}(\Env)}
     {\Env\vdash \texttt{catalyst}\ c:e\ \texttt{independent}\ \iota:\texttt{Catalyst}}
\end{align*}

\subsection{The two non-standard rules}

\begin{definition}[Positivity rule]
\label{def:positivity}
\[
\textsc{T-Seek-Pos}\quad
\frac{\Env\vdash t:\texttt{Cell}\quad \Env\vdash x:\texttt{Cell}\quad
      \Env\vdash \flo:\texttt{Floor}\quad \flo>0}
     {\Env\vdash \seek\ t\ \excl\ x\ \dots:\texttt{Outcome}}
\]
A \texttt{seek} is well typed only if the substrate's floor is positive. A
program asserting an attainable zero residual---by declaring a floor of
zero, or by requesting a point-valued target---does not type.
\end{definition}

\begin{definition}[Coherence rule]
\label{def:coherence-rule}
\[
\textsc{T-Seek-Coh}\quad
\frac{\Env\vdash c_i:\texttt{Catalyst}\ (1\le i\le k)\quad
      k\ge3\quad \mathrm{MutIndep}(c_1,\dots,c_k)}
     {\Env\vdash \seek\ \dots\ \via\ (c_1,\dots,c_k)\ \until\ \texttt{closure}:\texttt{Outcome}}
\]
where $\mathrm{MutIndep}$ holds when the declared independence relation
contains every unordered pair. A \texttt{seek} with an explicit \texttt{via}
clause is well typed only if it names at least three mutually independent
catalysts.
\end{definition}

\Cref{sec:coherence} justifies the threshold of three.

\subsection{Necessity of the exclusion clause}

\begin{theorem}[Exclusion cannot be eliminated]
\label{thm:exclusion-necessary}
Let a target be specified by a positive description alone---a predicate
asserting properties the target has---within a substrate whose state space
contains no distinguished element. Then the description does not determine a
unique cell, unless the description already names the complement it is to be
distinguished from. Consequently a \texttt{seek} without an exclusion clause
has no unique denotation, and the language is right to exclude the form.
\end{theorem}

\begin{proof}
A cell is an equivalence class of outcomes (\Cref{def:cell}); to determine a
cell is to determine, of each outcome, whether it belongs. A positive
description asserts properties of members. Suppose two distinct cells
$A\ne A'$ both satisfy every asserted property---which is possible whenever
the description does not mention a property on which $A$ and $A'$ differ.
Then the description determines neither, and no appeal to a distinguished
element resolves the tie, since by hypothesis none exists. To force
uniqueness the description must assert a property that $A$ has and $A'$
lacks; asserting such a property is asserting what the target is
\emph{not}---namely, not the thing lacking it---which is an exclusion. Hence
either the specification is incomplete or it contains an exclusion, and a
form with no exclusion clause cannot in general denote uniquely.
\end{proof}

\begin{corollary}[Compile-time rejection]
\label{cor:reject}
A parser for the grammar of \Cref{sec:grammar} rejects
$\seek\ t\ \via\ \dots$ without \texttt{excluding}, and no typing derivation
exists for such a term.
\end{corollary}

% =====================================================================
\section{Why three catalysts}
\label{sec:coherence}
% =====================================================================

\begin{definition}[Support structure]
\label{def:support}
Given catalysts $c_1,\dots,c_k$ bearing on a target, the \emph{support
structure} is the directed graph on $\{c_1,\dots,c_k\}$ with an edge
$c_i\to c_j$ when the outcome of $c_i$ is used to license the outcome of
$c_j$. A conclusion is \emph{robust} if it survives the removal of any
single catalyst.
\end{definition}

\begin{theorem}[Linear support is not robust]
\label{thm:linear}
If the support structure is acyclic, the conclusion is not robust: there is a
catalyst whose removal leaves some other catalyst's contribution
unlicensed, and by induction the chain collapses.
\end{theorem}

\begin{proof}
An acyclic finite digraph has a vertex of in-degree zero, say $c_s$: its
contribution is licensed by nothing internal to the structure. Removing the
terminal vertex of a maximal path leaves its predecessor's contribution
unused, and removing $c_s$ leaves every vertex reachable only through it
unlicensed. In either case some contribution's licensing is lost, so the
conclusion does not survive the removal of that single catalyst, and is not
robust by \Cref{def:support}.
\end{proof}

\begin{theorem}[Three is the minimum]
\label{thm:three}
A robust support structure contains a directed cycle of length at least
three. Cycles of length one \textup{(}self-support\textup{)} and two
\textup{(}mutual support of a pair\textup{)} do not confer robustness.
Conversely, three catalysts each supported by the other two form a robust
structure.
\end{theorem}

\begin{proof}
By \Cref{thm:linear} a robust structure is not acyclic, so it contains a
cycle. A $1$-cycle is a catalyst citing itself, which licenses nothing not
already assumed. For a $2$-cycle $\{c_1,c_2\}$: removing either leaves the
other supported by nothing, so the conclusion does not survive a single
removal and is not robust. Hence the shortest cycle compatible with
robustness has length $\ge3$. Conversely, in a $3$-cycle with all mutual
edges, removing any one catalyst leaves the remaining two still mutually
supporting via the edge between them, so a licensed contribution survives
and the structure is robust.
\end{proof}

\begin{remark}[What the rule does and does not guarantee]
\label{rem:coh-limits}
\Cref{def:coherence-rule} enforces a structural necessary condition, not
sufficiency for a correct conclusion. Three catalysts that are declared
independent but are not---sharing a common upstream data source, a common
preprocessing step, or a common modelling assumption---satisfy the rule
while violating its intent. The language cannot verify independence; it can
only require that it be declared, so that the declaration is visible in the
source and auditable. We regard this as the honest division of labour:
mechanical enforcement of the structure, human responsibility for the
premise.
\end{remark}

% =====================================================================
\section{Operational semantics}
\label{sec:semantics}
% =====================================================================

\subsection{Configurations}

\begin{definition}[Configuration]
\label{def:config}
A \emph{configuration} is a triple
$\langle e,\ \store,\ \rho\rangle$ where $e$ is the term being evaluated,
$\store$ is the value store mapping identifiers to values, and
$\rho\in\N$ is the \emph{committed record}, a count of catalyst invocations
performed.
\end{definition}

\begin{definition}[Reachable-cell set]
\label{def:reachable}
For a \texttt{seek} in progress, $\mathcal{C}$ denotes the set of cells
reached by catalysts invoked so far.
\end{definition}

\subsection{Reduction rules}

We write $\store(c)\Downarrow z$ for ``invoking catalyst $c$ against the
substrate yields cell $z$''. The rules are:

\begin{align*}
&\textsc{E-Invoke}\quad
\frac{c\in \Gamma_{\mathrm{avail}}\qquad \store(c)\Downarrow z}
{\langle \seek\ t\ \excl\ x\ \via\ \Gamma_{\mathrm{avail}}\ \until\ \texttt{closure},\ \store,\ \rho\rangle
\steps
\langle \seek\ t\ \excl\ x\ \via\ \Gamma_{\mathrm{avail}}\setminus\{c\}\ \until\ \texttt{closure},\ \store[\mathcal{C}\mapsto\mathcal{C}\cup\{z\}],\ \rho+1\rangle}\\[8pt]
&\textsc{E-Close-Res}\quad
\frac{\forall c\in\Gamma_{\mathrm{avail}}:\ \store(c)\Downarrow z\Rightarrow z\in\mathcal{C}
\qquad \abs{\mathcal{C}}=1}
{\langle \seek\ \dots,\ \store,\ \rho\rangle\steps
\langle \texttt{Resolved}(z_0),\ \store,\ \rho\rangle}\\[8pt]
&\textsc{E-Close-Dec}\quad
\frac{\forall c\in\Gamma_{\mathrm{avail}}:\ \store(c)\Downarrow z\Rightarrow z\in\mathcal{C}
\qquad \abs{\mathcal{C}}\ge2}
{\langle \seek\ \dots,\ \store,\ \rho\rangle\steps
\langle \texttt{Declined}(\mathcal{C}),\ \store,\ \rho\rangle}
\end{align*}

\textsc{E-Invoke} consumes one available catalyst, adds its cell to the
reached set, and increments the record. The two closure rules fire when no
remaining catalyst can add a cell, and are distinguished solely by the
cardinality of $\mathcal{C}$.

\subsection{Metatheory}

\begin{theorem}[Progress]
\label{thm:progress}
If $\Env\vdash e:\texttt{Outcome}$ and $e$ is not of the form
$\texttt{Resolved}(\cdot)$ or $\texttt{Declined}(\cdot)$, then for any store
$\store$ and record $\rho$ there is a configuration to which
$\langle e,\store,\rho\rangle$ steps.
\end{theorem}

\begin{proof}
A well-typed non-value term of type \texttt{Outcome} is a \texttt{seek} by
\Cref{def:positivity,def:coherence-rule}, the only rules concluding that
type. Consider its available set $\Gamma_{\mathrm{avail}}$. If some
available catalyst yields a cell outside $\mathcal{C}$, then
\textsc{E-Invoke} applies with that catalyst. Otherwise every available
catalyst yields a cell in $\mathcal{C}$; then $\mathcal{C}$ is nonempty---by
\textsc{T-Seek-Coh} at least three catalysts were named and, if none has yet
been invoked, \textsc{E-Invoke} applies to any of them---so exactly one of
\textsc{E-Close-Res} ($\abs{\mathcal{C}}=1$) or \textsc{E-Close-Dec}
($\abs{\mathcal{C}}\ge2$) applies. In every case a rule fires, so the term
is not stuck.
\end{proof}

\begin{theorem}[Preservation]
\label{thm:preservation}
If $\Env\vdash e:T$ and
$\langle e,\store,\rho\rangle\steps\langle e',\store',\rho'\rangle$ then
$\Env\vdash e':T$.
\end{theorem}

\begin{proof}
By cases on the rule applied. \textsc{E-Invoke} yields a \texttt{seek} term
differing from its predecessor only in the available set and the store; the
premises of \textsc{T-Seek-Pos} concern the floor, unchanged, and those of
\textsc{T-Seek-Coh} concern the catalysts named in the \texttt{via} clause
at the time of typing, which are not re-checked at each step---the rule
constrains the program text, not the residual available set. Hence the
reduct types at \texttt{Outcome}. \textsc{E-Close-Res} yields
$\texttt{Resolved}(z_0)$ with $z_0:\texttt{Cell}$, of type \texttt{Outcome}
by \Cref{def:outcome}; \textsc{E-Close-Dec} yields
$\texttt{Declined}(\mathcal{C})$ with $\abs{\mathcal{C}}\ge2$, likewise of
type \texttt{Outcome}. In all cases the type is preserved.
\end{proof}

\begin{corollary}[Soundness]
\label{cor:soundness}
A well-typed program does not get stuck: it either reduces indefinitely or
reaches a value of type \texttt{Outcome}. Combined with
\Cref{thm:termination} below, every well-typed program reaches such a value
in finitely many steps.
\end{corollary}

\begin{proof}
Immediate from \Cref{thm:progress,thm:preservation} by the standard
syntactic argument \citep{wright1994}.
\end{proof}

\begin{theorem}[Monotone record and non-return]
\label{thm:monotone}
The record $\rho$ is non-decreasing along any reduction and strictly
increases at each \textsc{E-Invoke}. Two configurations with identical terms
and stores but different records are distinct; hence no program returns to a
previous configuration.
\end{theorem}

\begin{proof}
\textsc{E-Invoke} increments $\rho$; the closure rules leave it unchanged
and produce values, from which no further step is taken. No rule decrements
$\rho$. Configurations with records $r<r'$ are separated by the predicate
$\rho\le r$, hence distinct. A return would require $\rho$ to decrease,
which no rule permits.
\end{proof}

\begin{theorem}[Determinism modulo substrate]
\label{thm:determinism}
Fix a substrate and fix, for each catalyst, the cell it yields. Then
reduction is deterministic up to the order of \textsc{E-Invoke} steps, and
the final \texttt{Outcome} is independent of that order.
\end{theorem}

\begin{proof}
The only choice in the semantics is which available catalyst
\textsc{E-Invoke} consumes. Each invocation adjoins $\store(c)\Downarrow z$
to $\mathcal{C}$, and set union is associative, commutative, and idempotent,
so the reached set after all catalysts have been consumed is
$\{\store(c)\Downarrow z : c\in\Gamma_{\mathrm{avail}}\}$ regardless of
order. The closure rules are selected by $\abs{\mathcal{C}}$, a function of
that set. Hence the final outcome is order-independent, and any residual
non-determinism arises solely from the substrate's responses.
\end{proof}

\begin{theorem}[Termination]
\label{thm:termination}
Every well-typed \texttt{seek} over a finite catalyst registry reaches an
\texttt{Outcome} in at most $\abs{\Gamma_{\mathrm{avail}}}$ invocations.
\end{theorem}

\begin{proof}
Each \textsc{E-Invoke} removes one catalyst from the available set, which is
finite, so at most $\abs{\Gamma_{\mathrm{avail}}}$ such steps occur. When the
set is empty the premise of the closure rules---that every available
catalyst yields a cell already in $\mathcal{C}$---is vacuously satisfied, so
a closure rule fires and a value is produced. Hence termination in at most
$\abs{\Gamma_{\mathrm{avail}}}$ invocations, and possibly earlier if closure
is detected before exhaustion.
\end{proof}

\begin{theorem}[Closure is strictly stronger than a threshold]
\label{thm:closure-stronger}
For any threshold $\theta$ on attained uncertainty there is a substrate and
program in which a partial invocation attains uncertainty below $\theta$
while an uninvoked catalyst would reach a cell outside $\mathcal{C}$; the
threshold criterion terminates, the closure criterion does not.
\end{theorem}

\begin{proof}
Take a substrate with two catalysts $c_A,c_B$ reaching incompatible cells
$z_A\ne z_B$, with the uncertainty attained at $z_A$ below $\theta$.
Invoking $c_A$ alone satisfies the threshold. But $c_B$ is available and
$z_B\notin\mathcal{C}=\{z_A\}$, so the premise of both closure rules fails
and the closure criterion requires \textsc{E-Invoke} on $c_B$. Hence the two
criteria differ on this program, with closure requiring strictly more.
\end{proof}

\begin{theorem}[Dichotomy of outcomes]
\label{thm:dichotomy}
Every terminating \texttt{seek} produces either
$\texttt{Resolved}(z)$ or $\texttt{Declined}(\mathcal{C})$ with
$\abs{\mathcal{C}}\ge2$, and never both.
\end{theorem}

\begin{proof}
By \Cref{thm:termination} a closure rule fires. The two rules have premises
$\abs{\mathcal{C}}=1$ and $\abs{\mathcal{C}}\ge2$, which are mutually
exclusive and, since $\mathcal{C}$ is nonempty at closure (at least one
catalyst has been invoked by \Cref{thm:progress}), jointly exhaustive.
\end{proof}

\begin{corollary}[Declination is a first-class result]
\label{cor:decline}
A program whose evidence is genuinely conflicting terminates normally,
returning the plurality of incompatible cells. This is a value of type
\texttt{Outcome}, not an error, and \texttt{report} emits it as such.
\end{corollary}

% =====================================================================
\section{Substrate neutrality}
\label{sec:neutrality}
% =====================================================================

The language depends on a substrate only through the four obligations of
\Cref{def:substrate}. We exhibit three bindings of different character to
show the obligations are discharge\-able without altering the language.

\begin{example}[Oscillatory-coherence substrate]
\label{ex:osc}
Receivers are recordings. The observable is a coherence index in $[0,1]$
computed from multi-channel phase relationships, mapped to uncertainty by an
order-reversing affine map so that higher coherence is lower uncertainty.
Events are transitions between adjacent segments carrying different labels,
typed by the ordered label pair. The floor is estimated by the asymptotic
method of \Cref{sec:diagnostics}.
\end{example}

\begin{example}[Discrete-transition substrate]
\label{ex:discrete}
Receivers are experimental conditions. States are members of a small finite
class set with an associated occupancy and dwell statistic; uncertainty is
derived from occupancy. Events are class-to-class transitions, typed by the
ordered class pair. This binding is appropriate when the underlying record
is a transition matrix rather than a time series.
\end{example}

\begin{example}[Latency substrate]
\label{ex:latency}
Receivers are participants. States are trials labelled by condition, with
uncertainty an affine function of response latency. Events are the addition
of an information source to a condition, typed by the source added. This
binding yields a zero-free-parameter prediction: the power of a compound
condition should equal the composition of the powers of its constituents.
\end{example}

\begin{proposition}[Neutrality]
\label{prop:neutrality}
No rule of \Cref{sec:grammar,sec:typing,sec:semantics} refers to any feature
of a substrate beyond the four obligations. Consequently a program is
well-typed and reduces identically under any substrate discharging them.
\end{proposition}

\begin{proof}
Inspection of the rules: the grammar mentions substrate fields by keyword
only; \textsc{T-Sub} requires well-formedness of the four fields and nothing
else; \textsc{T-Seek-Pos} requires only $\flo>0$, supplied by (S4);
\textsc{T-Seek-Coh} concerns catalysts, not substrates; and the reduction
rules use the substrate only through $\store(c)\Downarrow z$, which is (S2)
and (S3) mediated by a catalyst. No rule inspects the internal structure of
a receiver, observable, event, or floor.
\end{proof}

% =====================================================================
\section{Diagnostics: when a binding is uninformative}
\label{sec:diagnostics}
% =====================================================================

A binding may satisfy every obligation and still be incapable of supporting
an inference. The language cannot detect this, but a diagnostic computed
from substrate data can, and we specify it so that a negative result is
attributed correctly.

\begin{definition}[Type separation]
\label{def:separation}
For a corpus of events with types, let $\mathrm{Var}_{\mathrm{b}}$ be the
variance across types of the per-type mean power and
$\mathrm{Var}_{\mathrm{w}}$ the mean within-type variance. The
\emph{separation} is
$\eta=\mathrm{Var}_{\mathrm{b}}/(\mathrm{Var}_{\mathrm{b}}+\mathrm{Var}_{\mathrm{w}})\in[0,1]$.
\end{definition}

\begin{proposition}[Uninformative bindings]
\label{prop:uninformative}
If $\eta\approx0$, the per-type means are indistinguishable, predictions
formed from them are constant at fixed cascade length, and no comparison
among composition rules has discriminating power. A negative result obtained
in this regime is evidence about the observable, not about the process.
\end{proposition}

\begin{proof}
If all type means coincide at $\bar\kap$, any prediction formed as a
function of the type means alone depends only on the number of events
composed; at fixed length it is constant, so its sample variance vanishes
and any correlation with measured outcomes is undefined. Competing rules
computed from the same means are likewise constant and cannot be separated.
\end{proof}

\begin{remark}[The floor diagnostic]
\label{rem:floor-diag}
The second diagnostic concerns obligation (S4). Two estimators are
distinguishable from data: the sample-minimum estimator returns the least
observed uncertainty and is positive whenever the sample is nonempty; the
asymptotic estimator fits the limiting value of separation cost as the
sample grows and can return a value statistically indistinguishable from
zero. Only the second can fail. A substrate binding should declare which it
uses, and a reader should treat a positivity claim from the former as
carrying no evidential weight. \Cref{sec:validation} tests both estimators
against generating processes with and without a genuine floor.
\end{remark}

% =====================================================================
\section{Numerical verification}
\label{sec:validation}
% =====================================================================

\subsection{Method}

We implement a lexer, parser, type checker, and evaluator for the language
as specified, and exercise them against constructed programs and synthetic
substrates. The following are checked.

\emph{Grammar and parsing.} Well-formed programs parse; programs omitting
the \texttt{excluding} clause are rejected at parse time
(\Cref{cor:reject}).

\emph{Typing.} Programs declaring a non-positive floor are rejected by
\textsc{T-Seek-Pos}; programs naming fewer than three catalysts, or three
that are not pairwise declared independent, are rejected by
\textsc{T-Seek-Coh}; conforming programs type.

\emph{Progress and preservation.} Over randomly generated well-typed
programs and substrates, every intermediate configuration is checked to be
either a value or reducible, and the type of each reduct is checked to be
\texttt{Outcome} (\Cref{thm:progress,thm:preservation}).

\emph{Determinism and termination.} Each program is evaluated under
randomised catalyst orderings; the final outcome is checked to be
independent of order (\Cref{thm:determinism}), and the invocation count is
checked against the bound of \Cref{thm:termination}.

\emph{Record monotonicity.} The record is checked to increase strictly at
each invocation and never to decrease (\Cref{thm:monotone}).

\emph{Dichotomy.} Substrates are constructed to force each branch: one in
which all catalysts agree, yielding \texttt{Resolved}; one in which two
catalysts reach incompatible cells, yielding \texttt{Declined}. We check
that the latter terminates normally rather than raising
(\Cref{cor:decline}).

\emph{Closure versus threshold.} The construction of
\Cref{thm:closure-stronger} is instantiated, and we verify that a
threshold rule stops after one catalyst while closure continues and reports
a declination.

\emph{Diagnostics.} Both floor estimators are applied to generating
processes with and without a floor, and the separation $\eta$ is computed
for substrates of varying type separation, verifying that the uninformative
regime is detected (\Cref{prop:uninformative}).

\subsection{What would count as a failure}

A parse acceptance of an exclusion-free \texttt{seek} would falsify
\Cref{cor:reject}; an order-dependent outcome would falsify
\Cref{thm:determinism}; a stuck configuration would falsify
\Cref{thm:progress}; a run exceeding the invocation bound would falsify
\Cref{thm:termination}; a raised exception on conflicting evidence would
falsify \Cref{cor:decline}. Results are recorded in machine-readable form
alongside this paper.

\subsection{Results}
\label{sec:results}

All five checks passed.

\emph{Parsing and typing.} The reference program of
\Cref{sec:grammar} parsed and type-checked. The same program with the
\texttt{excluding} clause deleted was rejected at parse time, as
\Cref{cor:reject} requires: the rejection occurs before typing, because the
form is absent from the grammar rather than ill-typed within it. Declaring a
floor of $0$ or of $-3$ was rejected by \textsc{T-Seek-Pos}; a
\texttt{via} clause naming two catalysts was rejected by
\textsc{T-Seek-Coh}; and a \texttt{via} clause naming three catalysts of
which one was not declared mutually independent of the others was also
rejected, showing the rule tests the independence relation and not merely
the count.

\emph{Semantics.} Over $500$ randomly generated substrates, each evaluated
under all $3!=6$ orderings of the catalyst set---$3000$ evaluations in
total---there were zero violations of progress or preservation, zero of
determinism, zero of the termination bound, zero of record monotonicity, and
zero of the outcome dichotomy. Every evaluation terminated in at most
$\abs{\Gamma_{\mathrm{avail}}}=3$ invocations, as \Cref{thm:termination}
requires, and every configuration reduced to a value of type
\texttt{Outcome}.

\emph{Both branches occur.} Of the $500$ substrates, $296$ terminated in
\texttt{Resolved} and $204$ in \texttt{Declined}. We record this because a
dichotomy theorem is compatible with one branch being unreachable in
practice, and a language whose declination branch never fires would be a
language with a threshold criterion in disguise. The near-balance here is an
artefact of the substrate sampling---cell counts were drawn so that agreement
and disagreement were both common---but it establishes that both branches
are live and that declination arises from ordinary substrate configurations
rather than from contrived ones.

\emph{Closure against a threshold.} On the substrate in which two catalysts
reach a common cell and the third reaches an incompatible one, with the
common cell carrying an attained uncertainty of $5.0$ against a threshold of
$\theta=10.0$, the threshold procedure halted after $1$ invocation reporting
$\texttt{Resolved}(\mathtt{cellA})$, while the closure procedure performed
$2$ invocations and reported
$\texttt{Declined}(\{\mathtt{cellA},\mathtt{cellB}\})$. The two procedures
returned incompatible verdicts on identical evidence, which is the concrete
content of \Cref{thm:closure-stronger}.

\begin{remark}[The threshold procedure stopped before the disagreement was
reachable]
\label{rem:threshold-early}
The detail worth drawing out is \emph{when} the threshold procedure halted:
after one invocation, not two. It stopped as soon as a single source reported
an outcome it found confident, and therefore never invoked the source that
would have contradicted it---not because that source was unavailable, but
because the criterion had already been satisfied. This is the failure mode of
\Cref{rem:closure-why} realised in a concrete evaluation: internal
consistency of one line of evidence is sufficient to satisfy a threshold, and
a threshold cannot express the requirement that the remaining sources be
consulted. Closure expresses exactly that requirement, and does so without
reference to any magnitude---the closure rules of \Cref{sec:semantics} inspect
only whether the reached set would grow, never how confident any outcome is.
\end{remark}

\emph{Diagnostics.} The separation statistic distinguished an informative
from an uninformative binding: $\eta=0.992$ for a substrate whose four event
types had well-spaced mean powers, against $\eta=2.20\times10^{-3}$ for a
substrate whose type means agreed to three significant figures, the latter
falling below the flagging threshold of \Cref{prop:uninformative}.

The floor estimators were compared in two sampling regimes,
\Cref{tab:floor}. In both, the sample minimum returned a strictly positive
value for a process with no floor---$2.82\times10^{-11}$ at realistic sample
sizes and $6.08\times10^{-17}$ at large ones---and was biased upward on the
floored process. The asymptotic estimator recovered the true floor to within
$0.002$ and returned a negative value on the unfloored process in both
regimes. The contrast is the practical justification for requiring a
substrate to declare which estimator discharges obligation~(S4): the two
agree on a process that has a floor and disagree in kind on one that does
not, and only the second can report the disagreement.

\begin{table}[H]
\centering
\caption{Floor estimators on generating processes with true floor $10$ and
with no floor.}
\label{tab:floor}
\begin{tabular}{@{}llcc@{}}
\toprule
Regime & Estimator & Floored ($\flo=10$) & Unfloored ($\flo=0$)\\
\midrule
\multirow{2}{*}{$n=50$--$400$}
 & Sample minimum & $10.0000000010$ & $2.82\times10^{-11}$\\
 & Asymptotic     & $9.9981$        & $-1.76\times10^{-5}$\\
\midrule
\multirow{2}{*}{$n=50$--$4000$}
 & Sample minimum & $10.0000000000$ & $6.08\times10^{-17}$\\
 & Asymptotic     & $10.0000$       & $-2.32\times10^{-3}$\\
\bottomrule
\end{tabular}
\end{table}

% =====================================================================
\section{Scope and limitations}
\label{sec:scope}
% =====================================================================

The language expresses the structure of an inquiry and delegates all
computation to substrates; it is therefore useless without a substrate and
makes no claim to compute anything itself. The coherence rule enforces a
structural condition and cannot verify the independence it requires
(\Cref{rem:coh-limits}). The closure criterion is only as good as the
catalyst registry: closure over an impoverished registry is easily attained
and means only that the registry is exhausted, not that the question is
settled---the criterion quantifies over available catalysts, not over
conceivable ones. The positivity rule prevents a program from asserting a
zero residual but cannot prevent a substrate from discharging the floor
obligation trivially; the diagnostic of \Cref{rem:floor-diag} is offered
because the type system cannot do this work. Finally, the semantics treats
catalyst invocation as returning a determinate cell; substrates whose
responses are noisy require the indistinguishability relation of
\Cref{def:cell} to absorb that noise, and the choice of that relation is a
substantive modelling decision the language does not make.

\subsection{Conclusion}

Mekaneck has one primitive, two non-standard typing rules, and three
reduction rules. The primitive requires an inquiry to say what its target is
being told apart from, and the type system rejects programs that do not; the
termination condition requires an inquiry to exhaust its available means
rather than to cross a threshold, and the semantics guarantees it does so in
finitely many steps; and the outcome type admits an explicit declination, so
that conflicting evidence terminates a program normally instead of forcing a
choice the evidence does not support. The metatheory is standard---progress,
preservation, determinism modulo substrate, termination---and the substrate
interface is narrow enough that the same program runs unchanged over
substrates of quite different character. What the language cannot do, it
declines to pretend to do: it cannot verify independence, cannot validate a
floor estimator, and cannot know whether the catalysts available to it are
the ones that matter.

\bibliographystyle{plainnat}
\bibliography{references}

\end{document}
